% !TEX program = xelatex
% =====================================================================
%  EJEMPLO — UBCDSlidesTemplate
%  Muestra todos los componentes del tema. Copia este fichero, borra lo
%  que no uses y sustituye el contenido.
%  Compila con XeLaTeX:  latexmk main.tex      (ver latexmkrc)
%  Notas de orador: descomenta \setbeameroption{show notes on second screen=right}
% =====================================================================
\documentclass[aspectratio=169,t]{beamer}
\usetheme{UBCDSlidesTemplate}
\usepackage[spanish,es-noshorthands,es-nolayout]{babel}
\usepackage{tabularx}

% \setbeameroption{show notes on second screen=right}

% Rutas de imágenes: logos del tema, logos propios y figuras
\graphicspath{{../logos/}{logos/}{figures/}}

% ---------------- DATOS DE PORTADA ------------------------------------
\title{Título de la presentación}
\subtitle{Subtítulo que explica de qué trata, en una o dos líneas}
\author{Nombre Apellidos}
\date{Lugar · 1 ene 2027}

% Fila superior de logos (opcional). Sin ella, la fecha va abajo a la derecha.
% \ubcdtoplogos{\includegraphics[height=14pt]{logo-a.pdf}\hspace{16pt}\includegraphics[height=14pt]{logo-b.png}}

% Pie de la portada (opcional): financiación, logos institucionales...
% \titlegraphic{\includegraphics[width=290pt]{institucionales.pdf}}

\ubcdkicker{UNIDAD · INSTITUTO · CENTRO}        % etiqueta junto al logo
\ubcdfootertext{Título corto · Unidad}           % pie de cada diapositiva
\ubcdtitlemeta{%                                  % filas de metadatos de la portada
  \ubcdmetarow{AUTOR}{Nombre Apellidos · Cargo}
  \ubcdmetarow{PROYECTO}{Código · Nombre del proyecto}
  \ubcdmetarow{EQUIPO}{Responsable 1 · Responsable 2}
}

\begin{document}

% ============================ PORTADA ================================
\begin{frame}[plain]
  \titlepage
  \note{Saludo y contexto en una frase. (≈1 min)}
\end{frame}

% ====================== SECCIÓN 1 · TARJETAS =========================
\ubcdsection{Sección uno}{Descripción breve que aparece en la portadilla granate de la sección.}

\begin{frame}{Tarjetas en columnas}{EYEBROW EN MAYÚSCULAS}
  \begin{columns}[T,onlytextwidth]
    \column{0.325\textwidth}
    \begin{ubcdcard}[equal height group=demo]
      \ubcdtag{ETIQUETA · 01}\ubcdcardtitle{Primera tarjeta}
      Texto de la tarjeta. Usa \texttt{equal height group} para igualar alturas.
    \end{ubcdcard}
    \column{0.325\textwidth}
    \begin{ubcdcard}[equal height group=demo]
      \ubcdtag{ETIQUETA · 02}\ubcdcardtitle{Segunda tarjeta}
      Otro texto, más corto.
    \end{ubcdcard}
    \column{0.325\textwidth}
    \begin{ubcdcard}[equal height group=demo, ubcdnuevo]
      \ubcdtag{ETIQUETA · 03 · NUEVO}\ubcdcardtitle{Tarjeta destacada}
      Borde discontinuo carmesí con la opción \texttt{ubcdnuevo}.
    \end{ubcdcard}
  \end{columns}
  \vfill
  \ubcdband{CLAVE}{Banda granate para el mensaje principal de la diapositiva.}
  \note{Una nota de orador por diapositiva. (≈45 s)}
\end{frame}

% Fondo "papel" (más claro) para alternar el ritmo visual
{\ubcdpaper
\begin{frame}{Filas numeradas}{FONDO PAPEL}
  \ubcdlead{Entradilla en gris: una frase que introduce la diapositiva.}
  \vspace{6pt}
  \begin{ubcdrows}
    \ubcdrow{01}{Primer punto}{Descripción del primer punto en una o dos líneas.}
    \ubcdrow{02}{Segundo punto}{Descripción del segundo punto.}
    \ubcdrow{03}{Tercer punto}{Descripción del tercer punto.}
  \end{ubcdrows}
  \note{Tres ideas, una por fila. (≈40 s)}
\end{frame}
}

% ====================== SECCIÓN 2 · DATOS ============================
\ubcdsection{Sección dos}{Cifras, tablas, flujos y capturas.}

\begin{frame}{Cifras y tabla}{DATOS}
  \begin{columns}[T,onlytextwidth]
    \column{0.24\textwidth}\ubcdstat[ubcdCrimson]{120}{NÚCLEOS}
    \column{0.24\textwidth}\ubcdstat{2 TB}{MEMORIA RAM}
    \column{0.24\textwidth}\ubcdstat{50 TB}{DISCO ÚTIL}
    \column{0.24\textwidth}\ubcdstat[ubcdCrimson]{GPU}{NUEVA EN 2027}
  \end{columns}
  \vspace{10pt}
  {\ubcdtablefont\renewcommand{\arraystretch}{1.6}\setlength{\tabcolsep}{4.3pt}%
  \begin{tabularx}{\textwidth}{@{\hspace{4.3pt}}p{0.1\textwidth}X p{0.2\textwidth}p{0.13\textwidth}}
    \ubcdthead \ubcdth{ID} & \ubcdth{Descripción} & \ubcdth{Área} & \ubcdth{Estado}\\
    \ubcdtalt A1 & Primera fila (fondo blanco) & Área uno & Hecho\\
    A2 & Segunda fila (fondo transparente) & Área dos & En curso\\
    \ubcdtalt A3 & Tercera fila & Área tres & Previsto\\
  \end{tabularx}}
  \note{Cifras clave y tabla resumen. (≈50 s)}
\end{frame}

{\ubcdpaper
\begin{frame}{Flujos de pasos}{PROCEDIMIENTO}
  {\ttfamily\bfseries\ubcdsizelabel\color{ubcdStoneOne}FLUJO ESTÁNDAR\par}\vspace{1.4pt}
  \ubcdflow{4}%
  \ubcdstep{01}{Entrada}\ubcdarrow
  \ubcdstep{02}{Registro}\ubcdarrow
  \ubcdstep{03}{Ejecución}\ubcdarrow
  \ubcdstep{04}{Entrega}\par
  \vspace{10pt}
  {\ttfamily\bfseries\ubcdsizelabel\color{ubcdCrimson}FLUJO NUEVO\par}\vspace{1.4pt}
  \ubcdflownuevo{3}%
  \ubcdstep{01}{Identificar}\ubcdarrow
  \ubcdstephl{02}{Paso destacado}\ubcdarrow
  \ubcdstep{03}{Desplegar}\par
  \note{Dos flujos: el estándar y el nuevo (borde discontinuo). (≈30 s)}
\end{frame}
}

\begin{frame}{Filas con píldora y captura}{RESULTADOS}
  \begin{minipage}[t]{262pt}
    \begin{ubcdrows}
      \ubcdtoolrow{Herramienta A}{Descripción de lo que hace y por qué importa.}{\ubcdpill[ok]{PÚBLICO}}
      \ubcdtoolrow{Herramienta B}{Otra descripción, algo más larga para ver cómo ajusta el texto.}{\ubcdpill[warn]{EN DESARROLLO}}
      \ubcdtoolrow{Herramienta C}{Pendiente de revisión.}{\ubcdpill[bad]{BLOQUEADO}}
      \ubcdtoolrow{Herramienta D}{Estado neutro.}{\ubcdpill{BORRADOR}}
    \end{ubcdrows}
  \end{minipage}\hfill
  \begin{minipage}[t]{119pt}
    \vspace{2pt}
    \ubcdshot{119pt}{110pt}{placeholder.png}
    \ubcdcaption{Pie de figura en monoespaciada.}
  \end{minipage}
  \note{Estados con píldoras de color y una captura a la derecha. (≈40 s)}
\end{frame}

% ============================ CIERRE =================================
\ubcdclosing{Gracias}{¿Preguntas?}{correo@ejemplo.org · sitio.web}

\end{document}
